\documentclass[conference]{IEEEtran}

\usepackage{amsmath}
\usepackage{booktabs}
\usepackage{graphicx}
\usepackage{siunitx}
\usepackage[hidelinks]{hyperref}

\title{Participant-Disjoint Evaluation of COVID-19 Respiratory-Audio Models with Validation-Only Selection}

\author{
\IEEEauthorblockN{Nishant Harkut and Ritu Tiwari}
\IEEEauthorblockA{Department of Information Technology\\
ABV-IIITM Gwalior, India}
}

\begin{document}
\maketitle

\begin{abstract}
Audio-based COVID-19 screening models report AUROCs of 0.90--0.97, but evaluation design affects estimates. We analyzed 2,088 Coswara participants under participant-disjoint evaluation with training-only feature ranking and validation-only model selection. The selected cough--speech fusion achieved test AUROC 0.895 (95\% CI 0.852--0.933). Speech alone achieved 0.888. The fusion gain was not statistically significant (DeLong $p=0.62$). Metadata alone achieved AUROC 0.964; shuffled-label retraining produced AUROC 0.503. Results show strong internal discrimination under rigorous evaluation but do not establish external validity or clinical readiness.
\end{abstract}

\begin{IEEEkeywords}
COVID-19, respiratory audio, evaluation design, participant-disjoint validation, multimodal fusion
\end{IEEEkeywords}

\section{Introduction}

Audio-based COVID-19 screening has achieved reported AUROCs of 0.90--0.97 \cite{laguarta2020cough,bhattacharya2023coswara,aytekin2024hst}. However, evaluation design substantially affects estimates: Han \emph{et al.} showed biased participant allocation can inflate AUROC from 0.71 to 0.90 \cite{han2022realistic}, and Coppock \emph{et al.} found matching reduces audio AUROC from 0.85 to 0.62 while symptoms alone outperformed audio \cite{coppock2024audio}.

We analyze Coswara data under participant-disjoint evaluation with validation-only model selection. Contributions: (1) training-only feature ranking from 10,140 candidates; (2) validation-only model and fusion selection; (3) test results with bootstrap uncertainty; (4) metadata and shuffled-label baselines. The study evaluates internal performance only; temporal and cross-dataset validation require separate protocols.

\section{Materials and Methods}

\subsection{Dataset and Preprocessing}

The indexed Coswara release contained 2,746 participants \cite{bhattacharya2023coswara}. After excluding 632 with unresolved labels and 26 failing quality checks, 2,088 participants entered the modeling cohort. Participants were assigned to train (N=1,460), validation (N=312), and test (N=316) partitions (70/15/15) with stratification on label, age, and gender. All recordings from one participant remained in one partition.

Audio was resampled to 16~kHz, silence-trimmed, and padded to fixed durations (4~s for cough, 5~s for breathing and speech). Quality screening required minimum duration, maximum silence ratio (0.70), and maximum clipping ratio (0.01).

\subsection{Feature Extraction and Selection}

Three descriptor banks were merged: openSMILE ComParE 2016 (6,373 features), INTERSPEECH 2010 functionals (1,586 features), and a project acoustic bank (2,177 features), yielding 10,140 candidates \cite{eyben2010opensmile,schuller2013compare}.

\textbf{Critical:} Feature ranking was computed \emph{only} on training data. LightGBM gain was normalized within modality and averaged across modalities. The top 800 features were frozen for validation and test without recomputing ranks.

\subsection{Model Training and Selection}

Four classifier families were evaluated: LightGBM, XGBoost, CatBoost, and sigmoid-calibrated RBF-SVC \cite{ke2017lightgbm,chen2016xgboost,prokhorenkova2018catboost}. SMOTE was applied only to training data \cite{chawla2002smote}. For each modality, validation AUROC selected either one classifier or the mean of top performers.

Three fusion rules were evaluated on \emph{validation} participants: uniform mean, validation-weighted mean, and logistic stacking. The primary rule was uniform averaging; validation AUROC selected only which modalities to include. Test metrics were excluded from all selection decisions.

\subsection{Baselines and Evaluation}

Metadata-only model used demographic and recording-protocol variables. Shuffled-label control established absence of leakage.

Primary endpoint was participant-level AUROC. Nonparametric bootstrap (1,000 resamples) yielded 95\% confidence intervals. Paired DeLong tests compared AUROCs \cite{delong1988auc}.

\section{Results}

\subsection{Validation-Only Selection}

Cough plus speech achieved validation AUROC 0.842, ranking first among predefined uniform-mean combinations. Validation AUROC selected LightGBM for speech and four-model mean for cough and breathing.

\subsection{Test Performance}

Table~\ref{tab:results} shows test results. Speech achieved AUROC 0.888 (95\% CI 0.842--0.930). The cough--speech mean achieved AUROC 0.895 (95\% CI 0.852--0.933) on 314 participants with both modalities available. The fusion gain over speech was 0.007 AUROC, not statistically significant (DeLong $p=0.62$).

\begin{table}[!t]
\caption{Test Results with 95\% Bootstrap Confidence Intervals}
\label{tab:results}
\centering
\footnotesize
\begin{tabular}{lrrr}
\toprule
System & $n$ & AUROC (95\% CI) & AUPRC \\
\midrule
Breathing, top-4 mean & 312 & 0.828 (0.775--0.877) & 0.734 \\
Cough, top-4 mean     & 316 & 0.862 (0.812--0.908) & 0.804 \\
Speech, LightGBM      & 314 & 0.888 (0.842--0.930) & 0.833 \\
\midrule
\textbf{Cough+speech mean} & \textbf{314} & \textbf{0.895 (0.852--0.933)} & \textbf{0.862} \\
\bottomrule
\end{tabular}
\end{table}

\subsection{Baselines}

Metadata-only model achieved AUROC 0.964 (95\% CI 0.938--0.984). Shuffled-label retraining produced AUROC 0.503.

\subsection{Literature Context}

Table~\ref{tab:context} compares with prior Coswara studies.

\begin{table}[!t]
\caption{Comparison with Prior Coswara Studies}
\label{tab:context}
\centering
\footnotesize
\begin{tabular}{llc}
\toprule
Study & Protocol & AUROC \\
\midrule
Bhattacharya \cite{bhattacharya2023coswara} & Internal, 9 sounds + symptoms & 0.915 \\
Chetupalli \cite{chetupalli2023multimodal} & Subject-disjoint, audio fusion & 0.880 \\
Truong \cite{truong2024fair} & Fixed test, 7 sounds & 0.866 \\
\midrule
\textbf{This work} & Participant-disjoint, 2 sounds & \textbf{0.895} \\
\bottomrule
\end{tabular}
\par\vspace{1pt}\scriptsize \textbf{Note:} Not head-to-head comparisons.
\end{table}

\section{Discussion}

A multimodal audio model achieved AUROC 0.895 under participant-disjoint evaluation, comparable to prior internal estimates (0.88--0.92). The fusion gain over speech (0.007) was not significant. Metadata alone achieved higher AUROC (0.964), confirming Han and Coppock's finding that non-audio variables are strongly predictive \cite{han2022realistic,coppock2024audio}. Shuffled-label AUROC of 0.503 establishes absence of gross leakage.

Limitations: First, labels were self-reported, not PCR-confirmed by authors. Second, validation served multiple selection roles; nested evaluation would be more conservative. Third, only 314 participants had complete predictions. Fourth, this evaluates internal performance only; temporal and cross-dataset validation require separate protocols.

\section{Conclusion}

Under participant-disjoint evaluation with validation-only selection, a Coswara audio model achieved AUROC 0.895 (95\% CI 0.852--0.933). Metadata outperformed audio (0.964). The estimate applies only to the test cohort and does not establish external validity or clinical readiness.

\bibliographystyle{IEEEtran}
\bibliography{references}

\end{document}